\documentclass[11pt,a4paper]{article}
\usepackage{fontspec,amsmath,amssymb,amsthm,microtype}
\setmainfont{lmroman10-regular.otf}[BoldFont=lmroman10-bold.otf,ItalicFont=lmroman10-italic.otf,BoldItalicFont=lmroman10-bolditalic.otf]
\usepackage{xcolor}
\usepackage[most]{tcolorbox}
\usepackage{mathrsfs}
\providecommand{\lt}{<}
\providecommand{\gt}{>}
\setmonofont{DejaVuSansMono.ttf}[BoldFont=DejaVuSansMono-Bold.ttf,ItalicFont=DejaVuSansMono-Oblique.ttf,BoldItalicFont=DejaVuSansMono-BoldOblique.ttf,Scale=0.86]
\newfontfamily\leanfallback{STIXTwoMath-Regular.otf}[Scale=0.86]
\definecolor{leanrule}{gray}{0.62}
\definecolor{leanback}{gray}{0.962}
\definecolor{leantitle}{gray}{0.22}
\definecolor{leanmuted}{gray}{0.38}
\newtcolorbox{leanbox}[1]{enhanced,breakable,colback=leanback,colframe=leanrule,
  boxrule=0.4pt,arc=0.6mm,left=2.4mm,right=2.4mm,top=1.8mm,bottom=1.8mm,
  before skip=0.8em,after skip=1em,fonttitle=\small\bfseries,coltitle=leantitle,
  colbacktitle=leanback,titlerule=0.3pt,title={#1}}
\newlength{\leancharwidth}
\newenvironment{leancode}{\par\smallskip\begingroup\ttfamily\footnotesize\raggedright
  \settowidth{\leancharwidth}{M}\setlength{\parskip}{0pt}\setlength{\parindent}{0pt}
  \frenchspacing}{\par\endgroup\smallskip}
\newcommand{\leanline}[2]{\hangindent=\dimexpr#1\leancharwidth+4\leancharwidth\relax
  \hangafter=1\noindent\hspace*{#1\leancharwidth}#2\par}
\newcommand{\leansource}[1]{\par{\footnotesize\color{leanmuted}#1}\par}
\newcommand{\leanlabel}[1]{\par\smallskip\noindent{\small\bfseries #1}\par}
\usepackage[margin=25mm]{geometry}
\usepackage[unicode,hidelinks]{hyperref}
\hypersetup{pdfauthor={LangevinLean Project}}
\setlength{\emergencystretch}{3em}
\setlength{\parskip}{0.4em}
\setlength{\parindent}{0pt}
\allowdisplaybreaks[2]
\raggedbottom
\title{\texorpdfstring{Transferring the constructed corrector to real observables}{Transferring the constructed corrector to real observables}}
\author{LangevinLean Project}
\date{}
\begin{document}\maketitle
Draft mathematical exposition; these notes may contain mistakes.


An isometric real embedding \(i:R\to H\) into a complex Hilbert space has orthogonal projection \(ii^*\) onto its closed real range. If a semigroup preserves this range, its actual integrated defect, finite corrector sums and strong limit preserve it too. Extraction transfers the same Poisson identity and constants to the real space.


\section{\texorpdfstring{The genuine closed real subspace}{The genuine closed real subspace}}


\textbf{Claim.} The real adjoint satisfies \(i^*i=I\), and \(\Pi=ii^*\) fixes exactly \(i(R)\). This image is closed.



\textbf{Proof.} The adjoint identity follows from preservation of the real inner product by an isometry. Hence \(\Pi(ix)=ix\). Conversely a fixed vector \(y\) equals \(i(i^*y)\), so it belongs to the image. The image is therefore the equalizer of the continuous maps \(\Pi\) and the identity, which is closed.

\begin{leanbox}{Lean: \texttt{instance\_\allowbreak{}at\_\allowbreak{}26}, \texttt{projection}, \texttt{adjoint\_\allowbreak{}embed} and 3 more}
\leansource{\texttt{instance\_\allowbreak{}at\_\allowbreak{}26} \quad WassersteinMixing/\allowbreak{}ComplexRealCorrector.lean, lines 26–26}
\begin{leancode}
\leanline{0}{local\ instance\ :\ InnerProductSpace\ ℝ\ H\ :=\ InnerProductSpace.rclikeToReal\ ℂ\ H}
\end{leancode}
\leansource{\texttt{projection} \quad WassersteinMixing/\allowbreak{}ComplexRealCorrector.lean, lines 29–30}
\begin{leancode}
\leanline{0}{def\ projection\ (i\ :\ R\ →ₗᵢ[ℝ]\ H)\ :\ H\ →L[ℝ]\ H\ :=}
\leanline{2}{i.toContinuousLinearMap\ ∘L\ i.toContinuousLinearMap.adjoint}
\end{leancode}
\leansource{\texttt{adjoint\_\allowbreak{}embed} \quad WassersteinMixing/\allowbreak{}ComplexRealCorrector.lean, lines 33–37}
\begin{leancode}
\leanline{0}{theorem\ adjoint\_embed\ (i\ :\ R\ →ₗᵢ[ℝ]\ H)\ (x\ :\ R)\ :}
\leanline{4}{i.toContinuousLinearMap.adjoint\ (i\ x)\ =\ x}
\end{leancode}
\leansource{\texttt{projection\_\allowbreak{}embed} \quad WassersteinMixing/\allowbreak{}ComplexRealCorrector.lean, lines 39–42}
\begin{leancode}
\leanline{0}{@[simp]\ theorem\ projection\_embed\ (i\ :\ R\ →ₗᵢ[ℝ]\ H)\ (x\ :\ R)\ :}
\leanline{4}{projection\ i\ (i\ x)\ =\ i\ x}
\end{leancode}
\leansource{\texttt{projection\_\allowbreak{}fixed\_\allowbreak{}iff} \quad WassersteinMixing/\allowbreak{}ComplexRealCorrector.lean, lines 45–51}
\begin{leancode}
\leanline{0}{theorem\ projection\_fixed\_iff\ (i\ :\ R\ →ₗᵢ[ℝ]\ H)\ (y\ :\ H)\ :}
\leanline{4}{projection\ i\ y\ =\ y\ ↔\ y\ ∈\ Set.range\ i}
\end{leancode}
\leansource{\texttt{isClosed\_\allowbreak{}range} \quad WassersteinMixing/\allowbreak{}ComplexRealCorrector.lean, lines 54–59}
\begin{leancode}
\leanline{0}{theorem\ isClosed\_range\ (i\ :\ R\ →ₗᵢ[ℝ]\ H)\ :\ IsClosed\ (Set.range\ i)}
\end{leancode}
\leanlabel{How the Lean proof follows the argument}
The local real inner product is the real part of the complex one. \texttt{projection} is the actual adjoint composition. \texttt{projection\_\allowbreak{}fixed\_\allowbreak{}iff} characterizes its fixed vectors, and \texttt{is\allowbreak{}Closed\_\allowbreak{}range} uses continuity of the two equalizer maps.
\end{leanbox}

\section{\texorpdfstring{Preservation passes through integrals, sums and limits}{Preservation passes through integrals, sums and limits}}


\textbf{Claim.} If \(P_ti(R)\subseteq i(R)\), then \(D_hi(R)\subseteq i(R)\), \(G_Ni(R)\subseteq i(R)\) and \(G_hi(R)\subseteq i(R)\).



\textbf{Proof.} Apply the projection to \(ix-P_tix\); it fixes each integrand. Continuity of the projection lets it pass through the Bochner integral. Induction on the finite-sum recursion gives the same property for every \(G_N\). Taking the strong limit preserves the projection equality, so the constructed corrector also lies in the genuine closed real image.

\begin{leanbox}{Lean: \texttt{PreservesReal}, \texttt{op\_\allowbreak{}fixed}, \texttt{defect\_\allowbreak{}fixed} and 2 more}
\leansource{\texttt{PreservesReal} \quad WassersteinMixing/\allowbreak{}ComplexRealCorrector.lean, lines 63–64}
\begin{leancode}
\leanline{0}{def\ PreservesReal\ (S\ :\ ReversibleSemigroup\ H)\ (i\ :\ R\ →ₗᵢ[ℝ]\ H)\ :\ Prop\ :=}
\leanline{2}{∀\ t,\ 0\ ≤\ t\ →\ ∀\ x,\ projection\ i\ (S.op\ t\ (i\ x))\ =\ S.op\ t\ (i\ x)}
\end{leancode}
\leansource{\texttt{op\_\allowbreak{}fixed} \quad WassersteinMixing/\allowbreak{}ComplexRealCorrector.lean, lines 68–74}
\begin{leancode}
\leanline{0}{theorem\ op\_fixed\ (hreal\ :\ PreservesReal\ S\ i)\ \{t\ :\ ℝ\}\ (ht\ :\ 0\ ≤\ t)}
\leanline{4}{(y\ :\ H)\ (hy\ :\ projection\ i\ y\ =\ y)\ :}
\leanline{4}{projection\ i\ (S.op\ t\ y)\ =\ S.op\ t\ y}
\end{leancode}
\leansource{\texttt{defect\_\allowbreak{}fixed} \quad WassersteinMixing/\allowbreak{}ComplexRealCorrector.lean, lines 77–85}
\begin{leancode}
\leanline{0}{theorem\ defect\_fixed\ (hreal\ :\ PreservesReal\ S\ i)\ (h\ :\ ℝ)\ (hh\ :\ 0\ ≤\ h)\ (x\ :\ R)\ :}
\leanline{4}{projection\ i\ (S.defect\ h\ (i\ x))\ =\ S.defect\ h\ (i\ x)}
\end{leancode}
\leansource{\texttt{partialCorrector\_\allowbreak{}fixed} \quad WassersteinMixing/\allowbreak{}ComplexRealCorrector.lean, lines 88–97}
\begin{leancode}
\leanline{0}{theorem\ partialCorrector\_fixed\ (hreal\ :\ PreservesReal\ S\ i)\ (h\ :\ ℝ)\ (hh\ :\ 0\ ≤\ h)}
\leanline{4}{(N\ :\ ℕ)\ (x\ :\ R)\ :}
\leanline{4}{projection\ i\ (S.partialCorrector\ h\ hh\ N\ (i\ x))\ =\ S.partialCorrector\ h\ hh\ N\ (i\ x)}
\end{leancode}
\leansource{\texttt{corrector\_\allowbreak{}fixed} \quad WassersteinMixing/\allowbreak{}ComplexRealCorrector.lean, lines 100–109}
\begin{leancode}
\leanline{0}{theorem\ corrector\_fixed\ (hreal\ :\ PreservesReal\ S\ i)\ (h\ :\ ℝ)\ (hh\ :\ 0\ ≤\ h)\ (x\ :\ R)\ :}
\leanline{4}{projection\ i\ (S.corrector\ h\ hh\ (i\ x))\ =\ S.corrector\ h\ hh\ (i\ x)}
\end{leancode}
\leanlabel{How the Lean proof follows the argument}
\texttt{Preserves\allowbreak{}Real} expresses only a semigroup range property. \texttt{defect\_\allowbreak{}fixed} commutes the projection through the justified integral. \texttt{partial\allowbreak{}Corrector\_\allowbreak{}fixed} uses the finite recursion, and \texttt{corrector\_\allowbreak{}fixed} uses uniqueness of limits. The actual diffusion proves the range property separately.
\end{leanbox}

\section{\texorpdfstring{The exact real corrector and its unchanged constants}{The exact real corrector and its unchanged constants}}


\textbf{Construction.} Define \(P_t^{\mathbb R}=i^*P_ti\) and \(G_h^{\mathbb R}=i^*G_hi\). Range preservation gives \(iP_t^{\mathbb R}=P_ti\) and \(iG_h^{\mathbb R}=G_hi\). Consequently

\[\|G_h^{\mathbb R}\|\le h,\qquad (I-P_h^{\mathbb R})G_h^{\mathbb R}x=\int_0^h(x-P_t^{\mathbb R}x)\,dt.\]

\textbf{Proof.} Apply the isometry to the extracted operators and use the fixed projection identities. Norm preservation transfers the bound without a dimension loss. Adjoint extraction commutes with the integral. Injectivity of the embedding transfers the Poisson equation; preservation of its real inner product transfers the full Dirichlet energy estimate with factor \(h^2\).

\begin{leanbox}{Lean: \texttt{realOp}, \texttt{realCorrector}, \texttt{embed\_\allowbreak{}realOp} and 8 more}
\leansource{\texttt{realOp} \quad WassersteinMixing/\allowbreak{}ComplexRealCorrector.lean, lines 112–113}
\begin{leancode}
\leanline{0}{def\ realOp\ (t\ :\ ℝ)\ :\ R\ →L[ℝ]\ R\ :=}
\leanline{2}{i.toContinuousLinearMap.adjoint\ ∘L\ (S.op\ t).restrictScalars\ ℝ\ ∘L\ i.toContinuousLinearMap}
\end{leancode}
\leansource{\texttt{realCorrector} \quad WassersteinMixing/\allowbreak{}ComplexRealCorrector.lean, lines 116–118}
\begin{leancode}
\leanline{0}{def\ realCorrector\ (h\ :\ ℝ)\ (hh\ :\ 0\ ≤\ h)\ :\ R\ →L[ℝ]\ R\ :=}
\leanline{2}{i.toContinuousLinearMap.adjoint\ ∘L\ (S.corrector\ h\ hh).restrictScalars\ ℝ\ ∘L}
\leanline{4}{i.toContinuousLinearMap}
\end{leancode}
\leansource{\texttt{embed\_\allowbreak{}realOp} \quad WassersteinMixing/\allowbreak{}ComplexRealCorrector.lean, lines 121–122}
\begin{leancode}
\leanline{0}{theorem\ embed\_realOp\ (hreal\ :\ PreservesReal\ S\ i)\ \{t\ :\ ℝ\}\ (ht\ :\ 0\ ≤\ t)\ (x\ :\ R)\ :}
\leanline{4}{i\ (realOp\ S\ i\ t\ x)\ =\ S.op\ t\ (i\ x)}
\end{leancode}
\leansource{\texttt{embed\_\allowbreak{}realCorrector} \quad WassersteinMixing/\allowbreak{}ComplexRealCorrector.lean, lines 125–127}
\begin{leancode}
\leanline{0}{theorem\ embed\_realCorrector\ (hreal\ :\ PreservesReal\ S\ i)\ (h\ :\ ℝ)\ (hh\ :\ 0\ ≤\ h)\ (x\ :\ R)\ :}
\leanline{4}{i\ (realCorrector\ S\ i\ h\ hh\ x)\ =\ S.corrector\ h\ hh\ (i\ x)}
\end{leancode}
\leansource{\texttt{realCorrector\_\allowbreak{}norm\_\allowbreak{}le} \quad WassersteinMixing/\allowbreak{}ComplexRealCorrector.lean, lines 130–136}
\begin{leancode}
\leanline{0}{theorem\ realCorrector\_norm\_le\ (hreal\ :\ PreservesReal\ S\ i)\ (h\ :\ ℝ)\ (hh\ :\ 0\ ≤\ h)\ (x\ :\ R)\ :}
\leanline{4}{‖realCorrector\ S\ i\ h\ hh\ x‖\ ≤\ h\ *\ ‖x‖}
\end{leancode}
\leansource{\texttt{realCorrector\_\allowbreak{}opNorm\_\allowbreak{}le} \quad WassersteinMixing/\allowbreak{}ComplexRealCorrector.lean, lines 139–142}
\begin{leancode}
\leanline{0}{theorem\ realCorrector\_opNorm\_le\ (hreal\ :\ PreservesReal\ S\ i)\ (h\ :\ ℝ)\ (hh\ :\ 0\ ≤\ h)\ :}
\leanline{4}{‖realCorrector\ S\ i\ h\ hh‖\ ≤\ h}
\end{leancode}
\leansource{\texttt{realDefect} \quad WassersteinMixing/\allowbreak{}ComplexRealCorrector.lean, lines 145–146}
\begin{leancode}
\leanline{0}{def\ realDefect\ (h\ :\ ℝ)\ (x\ :\ R)\ :\ R\ :=}
\leanline{2}{i.toContinuousLinearMap.adjoint\ (S.defect\ h\ (i\ x))}
\end{leancode}
\leansource{\texttt{realDefect\_\allowbreak{}eq\_\allowbreak{}integral} \quad WassersteinMixing/\allowbreak{}ComplexRealCorrector.lean, lines 149–157}
\begin{leancode}
\leanline{0}{theorem\ realDefect\_eq\_integral\ (h\ :\ ℝ)\ (hh\ :\ 0\ ≤\ h)\ (x\ :\ R)\ :}
\leanline{4}{realDefect\ S\ i\ h\ x\ =\ ∫\ t\ in\ 0..h,\ x\ -\ realOp\ S\ i\ t\ x}
\end{leancode}
\leansource{\texttt{embed\_\allowbreak{}realDefect} \quad WassersteinMixing/\allowbreak{}ComplexRealCorrector.lean, lines 160–161}
\begin{leancode}
\leanline{0}{theorem\ embed\_realDefect\ (hreal\ :\ PreservesReal\ S\ i)\ (h\ :\ ℝ)\ (hh\ :\ 0\ ≤\ h)\ (x\ :\ R)\ :}
\leanline{4}{i\ (realDefect\ S\ i\ h\ x)\ =\ S.defect\ h\ (i\ x)}
\end{leancode}
\leansource{\texttt{realCorrector\_\allowbreak{}poisson} \quad WassersteinMixing/\allowbreak{}ComplexRealCorrector.lean, lines 164–172}
\begin{leancode}
\leanline{0}{theorem\ realCorrector\_poisson\ (hreal\ :\ PreservesReal\ S\ i)\ (h\ :\ ℝ)\ (hh\ :\ 0\ ≤\ h)\ (x\ :\ R)\ :}
\leanline{4}{realCorrector\ S\ i\ h\ hh\ x\ -\ realOp\ S\ i\ h\ (realCorrector\ S\ i\ h\ hh\ x)\ =}
\leanline{6}{∫\ t\ in\ 0..h,\ x\ -\ realOp\ S\ i\ t\ x}
\end{leancode}
\leansource{\texttt{realCorrector\_\allowbreak{}energy\_\allowbreak{}le} \quad WassersteinMixing/\allowbreak{}ComplexRealCorrector.lean, lines 175–185}
\begin{leancode}
\leanline{0}{theorem\ realCorrector\_energy\_le\ (hreal\ :\ PreservesReal\ S\ i)\ (h\ :\ ℝ)\ (hh\ :\ 0\ ≤\ h)}
\leanline{4}{\{s\ :\ ℝ\}\ (hs\ :\ 0\ ≤\ s)\ (x\ :\ R)\ :}
\leanline{4}{inner\ ℝ\ (realCorrector\ S\ i\ h\ hh\ x)}
\leanline{6}{(realCorrector\ S\ i\ h\ hh\ x\ -\ realOp\ S\ i\ s\ (realCorrector\ S\ i\ h\ hh\ x))\ ≤}
\leanline{6}{h\ \textasciicircum{}\ 2\ *\ inner\ ℝ\ x\ (x\ -\ realOp\ S\ i\ s\ x)}
\end{leancode}
\leanlabel{How the Lean proof follows the argument}
\texttt{real\allowbreak{}Op}, \texttt{real\allowbreak{}Corrector} and \texttt{real\allowbreak{}Defect} are explicit adjoint extractions of constructed objects. The embedding lemmas identify them exactly. \texttt{real\allowbreak{}Corrector\_\allowbreak{}poisson} and \texttt{real\allowbreak{}Corrector\_\allowbreak{}energy\_\allowbreak{}le} transfer the already proved complex identities through the isometry.
\end{leanbox}

\section{\texorpdfstring{Actual coordinate embeddings on square-integrable observables}{Actual coordinate embeddings on square-integrable observables}}


\textbf{Construction.} A pointwise real linear isometry induces an isometry on \(L^2\) by composition. For a finite-dimensional real Hilbert space, use an orthonormal basis, take its real coordinates and include them in complex coordinate space.



\textbf{Proof.} The pointwise norm equality gives equality of the integral square norms, so the lifted map is a genuine isometry of equivalence classes. The coordinate representation is an isometry by the orthonormal basis theorem; real-to-complex inclusion preserves the sum of squared coordinate norms. Composing these maps and then lifting to \(L^2\) preserves every observable norm exactly.

\begin{leanbox}{Lean: \texttt{mapLp}, \texttt{mapLp\_\allowbreak{}coeFn}, \texttt{euclideanOfReal} and 3 more}
\leansource{\texttt{mapLp} \quad WassersteinMixing/\allowbreak{}ComplexRealCorrector.lean, lines 197–205}
\begin{leancode}
\leanline{0}{def\ mapLp\ (i\ :\ E\ →ₗᵢ[ℝ]\ F)\ (μ\ :\ Measure\ A)\ :\ Lp\ E\ 2\ μ\ →ₗᵢ[ℝ]\ Lp\ F\ 2\ μ\ where}
\leanline{2}{toLinearMap\ :=\ i.toContinuousLinearMap.compLpₗ\ 2\ μ}
\leanline{2}{norm\_map'\ f\ :=\ by}
\leanline{4}{change\ ‖i.toContinuousLinearMap.compLp\ f‖\ =\ ‖f‖}
\leanline{4}{rw\ [Lp.norm\_def,\ Lp.norm\_def]}
\leanline{4}{congr\ 1}
\leanline{4}{apply\ eLpNorm\_congr\_norm\_ae\ (Lp.aestronglyMeasurable\ \_)\ (Lp.aestronglyMeasurable\ \_)}
\leanline{4}{filter\_upwards\ [i.toContinuousLinearMap.coeFn\_compLp\ f]\ with\ a\ ha}
\leanline{4}{rw\ [ha,\ LinearIsometry.coe\_toContinuousLinearMap,\ i.norm\_map]}
\end{leancode}
\leansource{\texttt{mapLp\_\allowbreak{}coeFn} \quad WassersteinMixing/\allowbreak{}ComplexRealCorrector.lean, lines 208–212}
\begin{leancode}
\leanline{0}{theorem\ mapLp\_coeFn\ (i\ :\ E\ →ₗᵢ[ℝ]\ F)\ (μ\ :\ Measure\ A)\ (f\ :\ Lp\ E\ 2\ μ)\ :}
\leanline{4}{mapLp\ i\ μ\ f\ =ᵐ[μ]\ fun\ a\ ↦\ i\ (f\ a)}
\end{leancode}
\leansource{\texttt{euclideanOfReal} \quad WassersteinMixing/\allowbreak{}ComplexRealCorrector.lean, lines 215–229}
\begin{leancode}
\leanline{0}{def\ euclideanOfReal\ (ι\ :\ Type*)\ [Fintype\ ι]\ :}
\leanline{4}{EuclideanSpace\ ℝ\ ι\ →ₗᵢ[ℝ]\ EuclideanSpace\ ℂ\ ι\ where}
\leanline{2}{toFun\ x\ :=\ WithLp.toLp\ 2\ (fun\ j\ ↦\ (x\ j\ :\ ℂ))}
\leanline{2}{map\_add'\ x\ y\ :=\ by}
\leanline{4}{ext\ j}
\leanline{4}{simp}
\leanline{2}{map\_smul'\ c\ x\ :=\ by}
\leanline{4}{ext\ j}
\leanline{4}{simp}
\leanline{2}{norm\_map'\ x\ :=\ by}
\leanline{4}{rw\ [EuclideanSpace.norm\_eq,\ EuclideanSpace.norm\_eq]}
\leanline{4}{congr\ 1}
\leanline{4}{apply\ Finset.sum\_congr\ rfl}
\leanline{4}{intro\ j\ \_}
\leanline{4}{simp}
\end{leancode}
\leansource{\texttt{coordinate} \quad WassersteinMixing/\allowbreak{}ComplexRealCorrector.lean, lines 233–237}
\begin{leancode}
\leanline{0}{def\ coordinate\ (R\ :\ Type*)\ [NormedAddCommGroup\ R]\ [InnerProductSpace\ ℝ\ R]}
\leanline{4}{[FiniteDimensional\ ℝ\ R]\ :}
\leanline{4}{R\ →ₗᵢ[ℝ]\ EuclideanSpace\ ℂ\ (Fin\ (Module.finrank\ ℝ\ R))\ :=}
\leanline{2}{(euclideanOfReal\ (Fin\ (Module.finrank\ ℝ\ R))).comp}
\leanline{4}{(stdOrthonormalBasis\ ℝ\ R).repr.toLinearIsometry}
\end{leancode}
\leansource{\texttt{coordinateLp} \quad WassersteinMixing/\allowbreak{}ComplexRealCorrector.lean, lines 240–243}
\begin{leancode}
\leanline{0}{def\ coordinateLp\ (R\ :\ Type*)\ [NormedAddCommGroup\ R]\ [InnerProductSpace\ ℝ\ R]}
\leanline{4}{[FiniteDimensional\ ℝ\ R]\ (μ\ :\ Measure\ A)\ :}
\leanline{4}{Lp\ R\ 2\ μ\ →ₗᵢ[ℝ]\ Lp\ (EuclideanSpace\ ℂ\ (Fin\ (Module.finrank\ ℝ\ R)))\ 2\ μ\ :=}
\leanline{2}{mapLp\ (coordinate\ R)\ μ}
\end{leancode}
\leansource{\texttt{coordinateLp\_\allowbreak{}norm} \quad WassersteinMixing/\allowbreak{}ComplexRealCorrector.lean, lines 247–249}
\begin{leancode}
\leanline{0}{theorem\ coordinateLp\_norm\ (R\ :\ Type*)\ [NormedAddCommGroup\ R]\ [InnerProductSpace\ ℝ\ R]}
\leanline{4}{[FiniteDimensional\ ℝ\ R]\ (μ\ :\ Measure\ A)\ (f\ :\ Lp\ R\ 2\ μ)\ :}
\leanline{4}{‖coordinateLp\ R\ μ\ f‖\ =\ ‖f‖}
\end{leancode}
\leanlabel{How the Lean proof follows the argument}
\texttt{map\allowbreak{}Lp} uses actual continuous linear pointwise composition and proves its norm equality by almost-everywhere representative comparison. \texttt{euclidean\allowbreak{}Of\allowbreak{}Real}, \texttt{coordinate} and \texttt{coordinate\allowbreak{}Lp} construct the explicit finite-dimensional embedding; \texttt{coordinate\allowbreak{}Lp\_\allowbreak{}norm} records its exact norm equality.
\end{leanbox}

\end{document}
